\chapter{Conclusion \& Future Work}
\label{chapter:conclusion}

This thesis investigated the design, implementation, and evaluation of a real-time AI-based simultaneous interpretation system from a system-engineering perspective. Motivated by the limitations of traditional interpretation infrastructures and the growing feasibility of AI-driven speech and language technologies, the work focused on constructing a modular, provider-agnostic pipeline capable of operating under the strict latency and stability constraints imposed by live interpretation scenarios.

The proposed system integrates streaming ASR, NMT, and TTS synthesis within a unified architecture that emphasizes incremental processing, controlled segmentation, and asynchronous execution. Rather than optimizing individual models, the thesis deliberately adopted a system-level viewpoint, demonstrating how orchestration, buffering strategies, and stabilization mechanisms can enable reliable real-time behavior when leveraging existing cloud-based AI services. The resulting pipeline supports multi-language output, modular component substitution, and compatibility with professional audio environments.

An extensive evaluation framework was developed to assess performance beyond conventional offline accuracy metrics. End-to-end latency, component-level delays, and segmentation-related stability measures were analyzed under standardized real-time conditions. The results show that the system maintains latency within acceptable bounds for simultaneous interpretation across both baseline and alternative configurations, indicating architectural robustness with respect to component changes. Component-level analysis further revealed that recognition and segmentation dominate the latency budget, while translation and synthesis remain comparatively stable and bounded. These findings reinforced the importance of system-level design decisions in shaping real-time performance.

Ablation studies on segmentation parameters provided deeper insight into the trade-offs between responsiveness and output stability. The results demonstrate that intuitive parameter choices do not always yield the expected behavior: adaptive stability mechanisms can improve session-level latency despite increased local instability, and overly aggressive segmentation can paradoxically increase effective delay. These observations highlight the complexity of incremental processing and emphasize that segmentation strategies must be evaluated holistically within the full pipeline context.

Beyond software-level evaluation, this thesis assessed the feasibility of integrating AI-based interpretation pipelines into professional interpretation environments. The analysis showed that, when audio formats, routing concepts, and timing characteristics are carefully aligned with established workflows, AI-generated interpretation output can be injected into existing channel-based hardware infrastructures. While this study does not claim certification readiness, it demonstrates that such integration is technically plausible using current technologies and reasonable engineering assumptions, thereby narrowing the gap between research prototypes and professional deployment contexts.

\section*{Future Work}

Several directions for future research emerge from this work. First, the evaluation could be extended to more diverse and realistic speech conditions, including spontaneous speech, multiple speakers, and varying speaking styles. Such scenarios would allow further investigation of segmentation stability and latency behavior under increased variability and would provide a more comprehensive assessment of real-world performance.

Second, future work could incorporate perceptual and user-centered evaluation methodologies. While this thesis deliberately focused on system-level latency and stability, subjective listening studies and user experience assessments would offer valuable insight into how measured metrics translate into perceived interpretation quality. Combining quantitative system metrics with qualitative feedback would strengthen the evaluation of practical usability.

Third, the integration of adaptive control mechanisms represents a promising avenue for further exploration. Dynamic adjustment of segmentation parameters, buffering strategies, or speaking-rate compensation based on observed system state could improve robustness across a wider range of operating conditions. Such mechanisms would need to be designed carefully to avoid introducing additional instability, as highlighted by the ablation results.

Finally, deeper integration with professional interpretation hardware could be explored, including bidirectional interaction with consoles, redundancy mechanisms, and failover strategies suitable for mission-critical environments. Addressing certification requirements, compliance standards, and long-term reliability considerations would be necessary steps toward operational deployment.

In conclusion, this thesis demonstrates that AI-based simultaneous interpretation is not merely a question of model performance but a broader system-engineering challenge. By focusing on architecture, orchestration, and real-time constraints, the work provides a structured foundation for future research and development aimed at bringing AI-driven interpretation systems closer to practical, professional use.
